\subsection{From Natural Numbers to Rational Numbers}

Starting from $\N$, we can construct:
\begin{enumerate}
  \item The integers $\Z$, by adjoining an additive identity $0$ and additive inverses.
  \item The rational numbers
        \[
          \Q=\left\{\frac{n}{m} : n,m\in\Z,\ m\ne0\right\}.
        \]
\end{enumerate}
We assume the basic arithmetic properties of $\Z$ and $\Q$ have been established.
Our goal is to define $\R$. Why do we need $\R$?

\subsection{Algebraic Numbers}

\begin{definition}{Algebraic Number}{algebraic-number}
  A number $x$ is \textbf{algebraic} if it satisfies a polynomial equation
  \[
    c_nx^n+c_{n-1}x^{n-1}+\cdots+c_1x+c_0
    =\sum_{k=0}^{n}c_kx^k=0,
  \]
  where $c_k\in\Z$, $c_n\ne0$, and $n\ge1$.
\end{definition}

\begin{theorem}{Rational Numbers Are Algebraic}{rationals-are-algebraic}
  Every rational number is algebraic.
\end{theorem}

\begin{proof}
  Let $r=m/n\in\Q$, where $m,n\in\Z$ and $n\ne0$.
  Then $r$ satisfies $nx-m=0$, so it is algebraic by Definition~\ref{def:algebraic-number}.
\end{proof}

\begin{definition}{Divisibility}{divisibility}
  For $k,m\in\Z$, we say \textbf{$k$ divides $m$}, written $k\mid m$, if
  \[
    \exists n\in\Z\text{ such that }m=kn.
  \]
\end{definition}

\subsection{The Rational Zeros Theorem}

\begin{theorem}{Rational Zeros Theorem}{rational-zeros}
  Suppose $r\in\Q$ satisfies
  \[
    \sum_{k=0}^{n}c_k r^k=0,
    \qquad c_k\in\Z,\quad c_n\ne0,\quad n\ge1.
  \]
  Write $r=c/d$ in lowest terms, where $c,d\in\Z$, $d\ne0$, and $c,d$ have no common factor other than $\pm1$.
  Then
  \[
    c\mid c_0,\qquad d\mid c_n.
  \]
\end{theorem}

\begin{theorem}{Corollary: No Rational Square Root of Two}{no-rational-square-root-two}
  The equation $x^2-2=0$ has no rational solution.
\end{theorem}

\begin{proof}
  If $r=c/d\in\Q$ in lowest terms were a solution, Theorem~\ref{thm:rational-zeros} would give
  \[
    c\mid -2,\qquad d\mid 1.
  \]
  Thus $r$ would have to be $\pm1$ or $\pm2$, none of which is a solution.
\end{proof}

\begin{theorem}{Not Every Algebraic Number Is Rational}{algebraic-not-rational}
  The rational numbers form a proper subset of the algebraic numbers.
\end{theorem}

\begin{proof}
  Every rational number is algebraic by Theorem~\ref{thm:rationals-are-algebraic}.
  The number $\sqrt{2}$ is algebraic because it satisfies $x^2-2=0$, but it is not rational by Corollary~\ref{thm:no-rational-square-root-two}.
\end{proof}

\subsubsection{Proof of the Rational Zeros Theorem}
% To be continued in the next lecture.
